% Part: lambda-calculus
% Chapter: introduction
% Section: computable-lambda

\documentclass[../../../include/open-logic-section]{subfiles}

\begin{document}

\olfileid{lam}{int}{lrp}
\olsection{संगणनक्षम फलने \usetoken{S}{lambda definable} असतात}

\begin{thm}
\ollabel{thm:computable-lambda}
प्रत्येक संगणनक्षम आंशिक फलन !!{lambda definable} असते.
\end{thm}

\begin{proof}
%READERNOTE{SOL6-A111: क्लिनी प्रमाण रूपाने total primitive-recursive functions, total predicate चे लघुतमीकरण आणि search ला force करणारे शेवटचे उपयोजन पुरते आहे. या सिद्धतेत naive partial composition किंवा partial primitive-recursion construction गृहीत धरलेली नाही.}
प्रत्येक आंशिक संगणनक्षम फलन~$f$ साठी बंद प्रतिनिधी~$F$
शोधायचा आहे. क्लिनीच्या प्रमाण रूप प्रमेयानुसार, total आदिम
पुनरावर्ती अंकफलने दर्शवणे, total आदिम पुनरावर्ती zero-test
चे अपरिबद्ध लघुतमीकरण करणे, आणि त्यानंतर योग्य output
काढणे पुरेसे आहे. Total आदिम पुनरावर्ती फलनांसाठी आरंभीची
फलने, total संयोजन आणि total आदिम पुनरावर्तन यांच्या पुढील
constructions वापरू. Search चे मूल्य अपरिभाषित असताना
output extractor ते टाकून देऊ नये म्हणून शेवटचे उपयोजन
ते search आधी force करेल. ही शेवटची construction
लघुतमीकरणाच्या विभागात देऊन या प्रमेयाची सिद्धता पूर्ण करू.
सर्व आंशिक फलनांसाठी सरळ संयोजन किंवा सरळ recursion
योग्य असल्याचे आधीच गृहीत धरण्याची गरज नाही.
\end{proof}

उर्वरित पुरावा अधिक वाचनीय व्हावा म्हणून आपण अधिक प्रचलित
संकेतन वापरू. उदाहरणार्थ, $Mxyz$ ऐवजी $M(x, y, z)$ लिहू.
यातून काही अर्थ सूचित होत असला, तरी प्रकारविरहित लॅम्डा
कलनातील पदांना निश्चित स्थानसंख्या नसते, हे लक्षात ठेवा;
म्हणून त्याच पदासाठी~$M$ आपण $M(x,y)$ आणि
$M(x,y,z,w)$ असे दोन्ही लिहू शकतो. पण या संकेतनाने
आपण $M$ ला तीन चरांचे फलन मानत आहोत, हे सूचित होते
आणि पुढील व्याख्यांमागील हेतू अधिक स्पष्ट होतो.
तसेच “$M$ ची व्याख्या $M =
\lambd[x][\lambd[y][\lambd[z][\dots]]]$ ने करा”
याऐवजी “$M$ ची व्याख्या $M(x,y,z) = \dots$ ने करा”
असे म्हणू.

\end{document}
